% GENERATED LOCAL-BUILD FLAT PART — DO NOT MERGE

% ---- BEGIN TIR/monograph/v12/chapters/ch04_c2_to_herm0_r3.tex ----
% !TEX root = ../../tir_monograph_v12.tex
\chapter[C2 to Herm0(2) and R3]{From $\mathbb C^2$ to $\operatorname{Herm}_0(2)\cong\mathbb R^3$}
\label{ch:v12_c2_herm0_r3}

\section{Traceless relational generator space}

Every Hermitian operator on the binary carrier has the Pauli decomposition
\begin{equation}
A=a_0I+\mathbf a\cdot\boldsymbol\sigma,
\qquad
\mathbf a=(a_x,a_y,a_z)\in\mathbb R^3.
\end{equation}
The uniform component $a_0I$ acts equally on the two primitive alternatives. The distinction-generating translation space is therefore
\begin{equation}
\boxed{
\mathfrak g_{\rm rel}
:=\operatorname{Herm}_0(2)
=\{A=A^\dagger:\operatorname{Tr}A=0\}
=\operatorname{span}_{\mathbb R}\{\sigma_x,\sigma_y,\sigma_z\}.
}
\label{eq:v12_herm0_def}
\end{equation}
Hence
\begin{equation}
\boxed{
\dim_{\mathbb R}\mathfrak g_{\rm rel}=3,
\qquad
\operatorname{Herm}_0(2)\cong\mathbb R^3.
}
\label{eq:v12_herm0_r3}
\end{equation}
\vTwelveStatus{A}{--}{PASS}
This is an exact linear-algebraic property of the admitted two-state complex carrier.

\section{Canonical metric}

Define the Hilbert--Schmidt pairing on traceless Hermitian generators by
\begin{equation}
\boxed{
\langle A,B\rangle_{\rm HS}
=\frac12\operatorname{Tr}(AB).
}
\label{eq:v12_hs_metric}
\end{equation}
The Pauli basis obeys
\begin{equation}
\frac12\operatorname{Tr}(\sigma_i\sigma_j)=\delta_{ij}.
\end{equation}
Thus, for
\[
A=\mathbf a\cdot\boldsymbol\sigma,
\qquad
B=\mathbf b\cdot\boldsymbol\sigma,
\]
one obtains
\begin{equation}
\boxed{
\langle A,B\rangle_{\rm HS}=\mathbf a\cdot\mathbf b,
\qquad
\|A\|^2=\frac12\operatorname{Tr}(A^2)=|\mathbf a|^2.
}
\label{eq:v12_hs_euclidean}
\end{equation}
The generator space therefore carries a positive-definite Euclidean quadratic form canonically inherited from the operator algebra.
\vTwelveStatus{A}{--}{PASS}

\section{Unit relational directions}

For
\begin{equation}
A=\mathbf n\cdot\boldsymbol\sigma,
\qquad
|\mathbf n|=1,
\end{equation}
the Pauli algebra gives
\begin{equation}
A^2=I.
\end{equation}
Conversely, a traceless Hermitian $2\times2$ operator with $A^2=I$ has unit coefficient vector. Hence
\begin{equation}
\boxed{
\{A\in\mathfrak g_{\rm rel}:\|A\|=1\}
\cong S^2.
}
\label{eq:v12_generator_sphere}
\end{equation}
The Bloch pure-state sphere and the unit generator sphere therefore live in the same coefficient carrier.

\section[SU(2) covariance]{$SU(2)$ covariance}

Conjugation by $U\in SU(2)$ acts as
\begin{equation}
A\mapsto UAU^\dagger.
\end{equation}
It preserves tracelessness, Eq.~\eqref{eq:v12_hs_metric}, and the norm. On coefficient vectors this is the adjoint rotation
\begin{equation}
\boxed{
\operatorname{Ad}:SU(2)\to SO(3),
\qquad
\ker\operatorname{Ad}=\{\pm I\}.
}
\end{equation}
Thus the same algebra supplies the three-real-dimensional carrier, its Euclidean metric and its rotation group.

\section{Trace-one affine state space}

Let a primitive locus $x$ carry a normalized two-level state
\begin{equation}
\rho_x=\rho_x^\dagger,
\qquad
\operatorname{Tr}\rho_x=1.
\end{equation}
The affine hull of normalized states is
\begin{equation}
\boxed{
\mathcal A_2
=\frac12I+\operatorname{Herm}_0(2),
\qquad
\dim_{\mathbb R}\mathcal A_2=3.
}
\label{eq:v12_affine_state_space}
\end{equation}
Its translation space is exactly $\operatorname{Herm}_0(2)$.

\section{Canonical affine relation}

For two primitive loci $x,y$, define
\begin{equation}
\boxed{
\mathcal E_{xy}:=2(\rho_y-\rho_x).
}
\label{eq:v12_affine_relation}
\end{equation}
If
\[
\rho_x=\frac12(I+\mathbf r_x\cdot\boldsymbol\sigma),
\qquad
\rho_y=\frac12(I+\mathbf r_y\cdot\boldsymbol\sigma),
\]
then
\begin{equation}
\boxed{
\mathcal E_{xy}
=(\mathbf r_y-\mathbf r_x)\cdot\boldsymbol\sigma
\in\operatorname{Herm}_0(2).
}
\end{equation}
The normalization factor $2$ removes the Bloch $1/2$ convention and gives the direct coefficient displacement.

The affine relation satisfies the exact endpoint laws
\begin{align}
\mathcal E_{xx}&=0,\\
\mathcal E_{yx}&=-\mathcal E_{xy},\\
\boxed{\mathcal E_{xz}}&=\boxed{\mathcal E_{xy}+\mathcal E_{yz}}.
\label{eq:v12_endpoint_composition}
\end{align}
Equivalently,
\begin{equation}
\mathcal E_{xy}+\mathcal E_{yz}+\mathcal E_{zx}=0.
\end{equation}
The local endpoint-composition rule is therefore inherited directly from affine state differences.
\vTwelveStatus{A}{--}{PASS}

\section{Metric of relational displacement}

Using Eq.~\eqref{eq:v12_hs_metric},
\begin{equation}
\boxed{
\|\mathcal E_{xy}\|^2
=\frac12\operatorname{Tr}(\mathcal E_{xy}^2)
=|\mathbf r_y-\mathbf r_x|^2.
}
\label{eq:v12_affine_relation_metric}
\end{equation}
A common local unitary frame transformation
\begin{equation}
\rho_x\mapsto U\rho_xU^\dagger
\end{equation}
gives
\begin{equation}
\mathcal E_{xy}\mapsto U\mathcal E_{xy}U^\dagger,
\end{equation}
so the relation, metric and composition law are all $SU(2)$-covariant with effective $SO(3)$ action on coefficient vectors.

\section{Spatial promotion gate}

The TIR spatial bridge is now sharply typed:
\begin{equation}
\boxed{
\texttt{RELATION\_AS\_CANONICAL\_QUANTUM\_STATE\_DIFFERENCE}:
\qquad
x\to y
\ \widehat{=}\ 
\mathcal E_{xy}=2(\rho_y-\rho_x).
}
\label{eq:v12_relation_promotion_rule}
\end{equation}
Under this identification, the local physical relation carrier is
\begin{equation}
\boxed{
T_x\Sigma\;\widehat{=}\;\operatorname{Herm}_0(2)\cong\mathbb R^3
}
\label{eq:v12_spatial_promotion}
\end{equation}
with metric inherited from Eq.~\eqref{eq:v12_hs_metric} and additive local displacement inherited from Eq.~\eqref{eq:v12_endpoint_composition}.
\vTwelveStatus{B}{--}{OPEN}
The exact algebraic consequences of the rule are established above; the foundational promotion of physical spatial relation to the canonical quantum-state difference remains the named bridge whose uniqueness is addressed by the next chapter.

\section{Tetrahedral affine frame preview}

Four affinely independent pure qubit states span the full trace-one affine hull. The regular tetrahedral set
\begin{equation}
\sum_{a=1}^4\mathbf n_a=0,
\qquad
\mathbf n_a\cdot\mathbf n_b=-\frac13\quad(a\neq b)
\end{equation}
provides a minimal four-point affine frame in this three-dimensional carrier. It later appears simultaneously as the regular spatial cell and the qubit informationally complete tetrahedral frame.

The output of Part I is therefore
\begin{equation}
\boxed{
\mathbb C^2
\to
\mathcal A_2=\frac12I+\operatorname{Herm}_0(2)
\to
\mathcal E_{xy}\in\operatorname{Herm}_0(2)\cong\mathbb R^3
}
\end{equation}
with exact affine composition, Euclidean metric and rotation covariance. Part II now asks when this carrier is uniquely promoted to local physical space and how its minimal finite cell closes.

% ---- END TIR/monograph/v12/chapters/ch04_c2_to_herm0_r3.tex ----
\clearpage

% ---- BEGIN TIR/monograph/v12/chapters/ch05_euclidean_geometry_dimension_gate.tex ----
% !TEX root = ../../tir_monograph_v12.tex
\chapter{Euclidean Geometry and the Spatial Dimension Gate}
\label{ch:v12_euclidean_dimension}

\section{Canonical affine relation carrier}

The normalized two-level state carrier has affine hull
\[
\mathcal A_2
=
\{\rho=\rho^\dagger:\operatorname{Tr}\rho=1\}
=
\frac12I+\operatorname{Herm}_0(2).
\]
Write
\[
V:=\operatorname{Herm}_0(2)
=
\operatorname{span}_{\mathbb R}\{\sigma_x,\sigma_y,\sigma_z\}.
\]
Since \(\dim_{\mathbb R}\operatorname{Herm}(2)=4\) and the trace-one condition removes one
real affine degree of freedom,
\[
\boxed{\dim_{\mathbb R}V=3.}
\]

For an ordered pair of admitted point states \((\rho_x,\rho_y)\), the affine torsor
has a unique endpoint-carrying displacement
\[
\boxed{\mathfrak e_{xy}:=\rho_y-\rho_x\in V},
\qquad
\rho_x+\mathfrak e_{xy}=\rho_y.
\]
The generator-normalized relation used below is
\[
\boxed{\mathcal E_{xy}:=2\mathfrak e_{xy}.}
\]
It satisfies the exact endpoint identities
\[
\mathcal E_{yx}=-\mathcal E_{xy},
\qquad
\boxed{\mathcal E_{xz}=\mathcal E_{xy}+\mathcal E_{yz}},
\]
and therefore
\[
\mathcal E_{xy}+\mathcal E_{yz}+\mathcal E_{zx}=0.
\]
This endpoint-carrying universal property is the canonical local relation exported
by the v12 spatial branch.

\section{Invariant Euclidean metric}

On \(V\), define
\[
\boxed{
\langle A,B\rangle
=
\frac12\operatorname{Tr}(AB).
}
\]
The Pauli basis obeys
\[
\frac12\operatorname{Tr}(\sigma_i\sigma_j)=\delta_{ij}.
\]
Hence, for
\[
A=\mathbf a\cdot\boldsymbol\sigma,
\qquad
B=\mathbf b\cdot\boldsymbol\sigma,
\]
one has
\[
\boxed{
\langle A,B\rangle=\mathbf a\cdot\mathbf b,
\qquad
\|A\|^2=\frac12\operatorname{Tr}(A^2)=|\mathbf a|^2.
}
\]
The positive quadratic form is therefore Euclidean in Pauli coefficients.

The Pythagorean relation follows immediately. If
\[
\langle A,B\rangle=0,
\]
then
\[
\begin{aligned}
\|A+B\|^2
&=\langle A+B,A+B\rangle\\
&=\|A\|^2+2\langle A,B\rangle+\|B\|^2,
\end{aligned}
\]
so
\[
\boxed{\|A+B\|^2=\|A\|^2+\|B\|^2.}
\]
Thus the local affine relation carrier already contains the Euclidean norm,
orthogonality and Pythagorean closure used by the later spatial construction.

\section{Rotational covariance}

For \(U\in SU(2)\),
\[
A\mapsto UAU^\dagger
\]
preserves trace and the Hilbert--Schmidt metric.  The induced action on coefficient
vectors is
\[
\boxed{
\operatorname{Ad}:SU(2)\longrightarrow SO(3),
\qquad
SU(2)/\{\pm I\}\cong SO(3).
}
\]
The unit relation directions therefore form
\[
\boxed{
\{A\in V:\|A\|=1\}\cong S^2.
}
\]

\section{Rank-three stabilization}

At a relational locus \(x\), define the outgoing relation span
\[
W_x
:=
\operatorname{span}_{\mathbb R}
\{\mathcal E_{xy}:y\sim x\}
\subseteq V.
\]
Assume a populated primitive patch,
\[
W_x\neq\{0\},
\]
and full local relational isotropy,
\[
R(W_x)=W_x
\qquad
\text{for every }R\in SO(3)
\]
from the adjoint action above.

The defining real representation of \(SO(3)\) on \(\mathbb R^3\) is irreducible.
Its invariant linear subspaces are therefore \(\{0\}\) and the full carrier.
Consequently
\[
\boxed{
W_x\neq0,\quad SO(3)W_x=W_x
\Longrightarrow
W_x=V
\Longrightarrow
\dim W_x=3.
}
\]
Rank one selects an axis and rank two selects a plane; the full primitive isotropy
stabilizes the three-dimensional carrier.

\section{Minimal faithful realization crosscheck}

The same dimension is independently fixed by the minimal faithful real orthogonal
realization of the primitive rotation group. A continuous \(SO(3)\) action on a
one-dimensional real orthogonal carrier has trivial connected image.  A
two-dimensional orthogonal image lies in \(SO(2)\), while \(SO(3)\) is non-Abelian
and perfect.  The defining three-dimensional representation is faithful.

Therefore the minimal faithful real orthogonal carrier has
\[
\boxed{\dim_{\mathbb R}V_x=3}
\]
and is unique up to orthogonal equivalence.  Its invariant positive metric is
unique up to an overall positive scale.  The canonical dimensionless
normalization used here is exactly
\[
h(A,B)=\frac12\operatorname{Tr}(AB).
\]

\section{Continuum interface}

The discrete spatial branch is now source-bound by
\[
\boxed{
(\rho_x,\rho_y)
\longmapsto
\mathcal E_{xy}=2(\rho_y-\rho_x)
\in\operatorname{Herm}_0(2)\cong\mathbb R^3.
}
\]
The remaining geometric completion gate is the regular gluing/refinement step
that turns neighboring local copies of \(V\) into a smooth coframe/tangent-bundle
carrier.  Chapter~\ref{ch:v12_connections_holonomy} supplies the discrete
connection, solder and torsion objects required for that refinement.

\section{Source and validation receipt}

The canonical theorem surfaces are
\begin{itemize}
\item \path{TIR/subrepos/the-space-of-geometry/foundations/CANONICAL_SPATIAL_RELATION_EXTRACTION_V0_1.md};
\item \path{TIR/foundations/TIR_RANK3_ISOTROPY_STABILIZATION_V0_1.md};
\item \path{TIR/subrepos/the-space-of-geometry/foundations/SPATIAL_PROMOTION_UNIQUENESS_V0_1.md}.
\end{itemize}

The deterministic v12 receipt is generated by
\path{TIR/validation/tir_v12_ch05_euclidean_spatial_gate_v0_1.py}.

% ---- END TIR/monograph/v12/chapters/ch05_euclidean_geometry_dimension_gate.tex ----
\clearpage

% ---- BEGIN TIR/monograph/v12/chapters/ch06_tetrahedral_closure.tex ----
% !TEX root = ../../tir_monograph_v12.tex
\chapter{Minimal Finite Cell and Tetrahedral Closure}
\label{ch:v12_tetrahedral_closure}

\section{Minimal finite isotropy}

Let
\[
\mathbf n_1,\ldots,\mathbf n_m\in\mathbb R^3,
\qquad
|\mathbf n_a|=1,
\]
be equal-weight local relation directions in the carrier constructed in
Chapter~\ref{ch:v12_euclidean_dimension}.  Define
\[
\mathbf M:=\sum_{a=1}^{m}\mathbf n_a,
\qquad
Q:=\sum_{a=1}^{m}\mathbf n_a\mathbf n_a^T.
\]
The finite isotropy conditions are
\[
\boxed{\mathbf M=0},
\qquad
\boxed{Q=\frac m3 I_3}.
\]
The second condition has rank three.  For \(m\le2\) the relation span is
rank-deficient.  If \(m=3\), the zero-moment equation places the third vector in
the span of the first two, so the rank is at most two.  Hence
\[
\boxed{m\ge4.}
\]

For \(m=4\), place the four vectors into
\[
N=(\mathbf n_1\ \mathbf n_2\ \mathbf n_3\ \mathbf n_4).
\]
Then
\[
NN^T=\frac43I_3,
\qquad
N\mathbf1=0.
\]
The Gram matrix \(G=N^TN\) is therefore
\[
\boxed{
G
=
\frac43I_4-\frac13\mathbf1\mathbf1^T,
}
\]
so
\[
\boxed{
\mathbf n_a\cdot\mathbf n_b
=
-\frac13
\quad(a\neq b).
}
\]
The minimal equal-weight finite full-isotropy cell is consequently the regular
tetrahedral frame.

A convenient exact representative is
\[
\mathbf n_1=\frac1{\sqrt3}(1,1,1),\quad
\mathbf n_2=\frac1{\sqrt3}(1,-1,-1),
\]
\[
\mathbf n_3=\frac1{\sqrt3}(-1,1,-1),\quad
\mathbf n_4=\frac1{\sqrt3}(-1,-1,1).
\]

\section{Exact shape data}

For distinct vertices,
\[
\|\mathbf n_a-\mathbf n_b\|^2
=
2-2\left(-\frac13\right)
=
\frac83,
\]
hence the normalized edge length is
\[
\boxed{\hat a=\sqrt{\frac83}}.
\]
The Euclidean volume of the regular tetrahedron is
\[
\boxed{
\hat V_{\Delta^3}
=
\frac{\hat a^3}{6\sqrt2}
=
\frac{8}{9\sqrt3}.
}
\]

\section{Independent qubit informational-completeness route}

A generic qubit density operator has three independent real Bloch coordinates.
An \(m\)-outcome normalized probability vector has at most \(m-1\) independent
real entries.  Informational completeness therefore requires
\[
m-1\ge3,
\qquad
\boxed{m\ge4}.
\]

Using the same tetrahedral directions, define
\[
P_a=\frac12(I+\mathbf n_a\cdot\boldsymbol\sigma),
\qquad
E_a=\frac12P_a
=\frac14(I+\mathbf n_a\cdot\boldsymbol\sigma).
\]
Since \(\sum_a\mathbf n_a=0\),
\[
\boxed{\sum_{a=1}^{4}E_a=I}.
\]
For \(a\neq b\),
\[
\boxed{\operatorname{Tr}(P_aP_b)=\frac13}.
\]
For
\[
\rho=\frac12(I+\mathbf r\cdot\boldsymbol\sigma)
\]
the outcome probabilities are
\[
p_a=\operatorname{Tr}(\rho E_a)
=\frac14(1+\mathbf r\cdot\mathbf n_a).
\]
Using
\[
\sum_a\mathbf n_a\mathbf n_a^T=\frac43I_3
\]
gives the exact reconstruction
\[
\boxed{
\mathbf r
=
3\sum_{a=1}^{4}p_a\mathbf n_a.
}
\]
Thus the minimal symmetric informationally complete qubit frame and the minimal
finite isotropic relation frame carry the same regular tetrahedral Gram data.

\section{Congruence-class closure}

Let \(N\) be an ordered spatial tetrahedral frame and \(M\) an ordered SIC
tetrahedral frame, with
\[
N^TN=M^TM=G,
\qquad
NN^T=MM^T=\frac43I_3.
\]
Define
\[
\boxed{Q_{\Delta}:=\frac34MN^T.}
\]
Then
\[
Q_{\Delta}N=M,
\qquad
Q_{\Delta}Q_{\Delta}^T=I_3,
\]
so
\[
\boxed{Q_{\Delta}\in O(3).}
\]
After compatible orientation of the ordered frames the representative may be
chosen in \(SO(3)\).  The two routes therefore occupy one exact tetrahedral
orthogonal-congruence class in
\(\operatorname{Herm}_0(2)\cong\mathbb R^3\).

The common promoted object at this stage is this Gram/congruence class and its
orthogonal shape invariants.  Semantic role binding between an informational
outcome label and a spatial adjacency label remains a separate downstream
binding.

\section{Fubini--Study shape crosswalk}

The Bloch central angle between tetrahedral directions obeys
\[
\cos\chi=-\frac13.
\]
For one spherical tetrahedral face the spherical cosine law gives the interior
angle
\[
\alpha=\frac{2\pi}{3}.
\]
The spherical excess on the unit Bloch sphere is \(\pi\).  Since the qubit
Fubini--Study metric satisfies
\[
ds_{FS}^2=\frac14ds_{S^2}^2,
\]
one face has
\[
\boxed{a_{FS}^{face}=\frac\pi4},
\]
and all four faces carry
\[
\boxed{a_{FS}^{tet}=\pi}.
\]
Therefore the dimensionless dual-shape coefficient is
\[
\boxed{
C_{\Delta/FS}
=
\frac{\hat V_{\Delta^3}}{a_{FS}^{tet}}
=
\frac{8}{9\sqrt3\,\pi}.
}
\]

\section{Legacy tetrahedral semantics}

The historical v11 tetrahedron appendix carries the sector-specific
operator labels
\(\{\hat I,\hat M,\hat A_\mu,\hat A_\tau\}\) and its own six edge assignments.
In v12 that object remains sector-level provenance.  The canonical geometric
tetrahedron of this chapter is the regular Gram class derived above; any
role-binding map between the two tetrahedral descriptions must preserve their
typed edge data and receives its own source receipt.

\section{Source and validation receipt}

The canonical theorem surfaces are
\begin{itemize}
\item \path{TIR/foundations/TIR_MINIMAL_ISOTROPIC_TETRAHEDRAL_CELL_V0_1.md};
\item \path{TIR/foundations/TIR_QUBIT_TETRAHEDRAL_INFORMATIONAL_COMPLETENESS_V0_1.md};
\item \path{TIR/integration/TIR_TETRAHEDRAL_CONGRUENCE_CLASS_CLOSURE_V0_1.md}.
\end{itemize}
The deterministic v12 receipt is generated by
\path{TIR/validation/tir_v12_ch06_tetrahedral_closure_v0_1.py}.

% ---- END TIR/monograph/v12/chapters/ch06_tetrahedral_closure.tex ----
\clearpage

% ---- BEGIN TIR/monograph/v12/chapters/ch07_connections_holonomy_se3_solder_torsion.tex ----
% !TEX root = ../../tir_monograph_v12.tex
\chapter[Connections, Holonomy, SE(3), Solder and Torsion]{Connections, $W_{ij}$, Holonomy, $SE(3)$, Solder and Torsion}
\label{ch:v12_connections_holonomy}

\paragraph{v12.1 synchronization.}
The original v12 migration closed this chapter at the discrete Universal-Loop
source and treated Cartan refinement as the next theorem.  The current
repository has since closed the local/refining Gates A2--A4 and implemented the
Gate-A5 global three-manifold certifier.  This chapter now records those results
without promoting the still-missing production global relational-complex input.

\section{Typed connection transport}

Let \(G\) be the structure group of an active sector and \(A_R\) its connection
in a representation \(R\).  The common transport object is
\[
\boxed{
W_{ij}^{(G,R)}
=\mathcal P\exp\!\left(\int_{\gamma_{ij}}A_R\right).
}
\]
For unitary transport,
\[
\boxed{W_{ji}=W_{ij}^{-1}=W_{ij}^{\dagger}},
\]
and under local frame changes
\[
\boxed{W_{ij}\mapsto G_iW_{ij}G_j^{-1}}.
\]

The v12 typing is
\[
\boxed{
W_{ij}^{WT}\in U(1),
\qquad W_{ij}^{X}\in SU(2),
\qquad W_{ij}^{c}\in SU(3).
}
\]
For the spatial branch,
\[
\boxed{R_{ij}:=\operatorname{Ad}_{W_{ij}^{X}}\in SO(3)}.
\]
A composable closed path carries
\[
\boxed{W_\gamma=W_{i_0i_1}W_{i_1i_2}\cdots W_{i_{m-1}i_0}}.
\]
Its conjugacy class is the gauge-covariant finite loop datum.

\section{Atlas transitions and the holonomy firewall}

Let a local affine chart \(a\) have anchor \(\mathbf r_a\) and orthonormal
frame \(Q_a\in SO(3)\), with
\[
x_a=Q_a^T(\mathbf r-\mathbf r_a).
\]
For overlapping charts,
\[
\boxed{x_b=Q_b^TQ_a\,x_a+Q_b^T(\mathbf r_a-\mathbf r_b)}.
\]
Therefore
\[
\boxed{
G_{ba}^{atlas}
=(R_{ba},t_{ba})
=\left(Q_b^TQ_a,Q_b^T(\mathbf r_a-\mathbf r_b)\right)
\in SE(3).
}
\]
For three charts,
\[
G_{cb}^{atlas}G_{ba}^{atlas}=G_{ca}^{atlas},
\qquad
G_{ab}^{atlas}=(G_{ba}^{atlas})^{-1}.
\]
A closed sequence of pure atlas transitions therefore gives
\[
\boxed{G_C^{atlas}=e_{SE(3)}}.
\]
This exact cocycle is the zero-holonomy baseline.  Path-dependent connection
transport is carried by the separately typed \(W_{ij}\) layer.

\section{Discrete solder object and endpoint defect}

For every oriented spatial edge define
\[
\mathcal E_{xy}=E^a_{xy}\sigma_a\in\operatorname{Herm}_0(2),
\qquad
\ell_{xy}^2=\frac12\operatorname{Tr}(\mathcal E_{xy}^2).
\]
The reversed edge obeys the covariant orientation rule
\[
\boxed{
\mathcal E_{yx}
=-W_{yx}^{X}\mathcal E_{xy}(W_{yx}^{X})^\dagger.
}
\]
For the oriented triangle \(x\to y\to z\to x\), define
\[
\boxed{
\mathcal T_{xyz}
=\mathcal E_{xy}
+W_{xy}^{X}\mathcal E_{yz}(W_{xy}^{X})^\dagger
+W_{xz}^{X}\mathcal E_{zx}(W_{xz}^{X})^\dagger.
}
\]
The transported endpoint defect is
\[
\boxed{
\mathcal C_{xyz}
=\mathcal E_{xz}
-\left(
\mathcal E_{xy}
+W_{xy}^{X}\mathcal E_{yz}(W_{xy}^{X})^\dagger
\right).
}
\]
The existing relational endpoint identity is
\[
\boxed{\mathcal T_{xyz}=-\mathcal C_{xyz}}.
\]

\section{Gate A: affine-loop torsion source}

Writing
\[
\mathcal E_{xy}=\mathbf e_{xy}\cdot\boldsymbol\sigma,
\qquad
R_{xy}=\operatorname{Ad}(W_{xy}^{X}),
\]
the connection-lifted affine edge is
\[
\boxed{G_{xy}^{\nabla}=(R_{xy},\mathbf e_{xy})\in SE(3)}.
\]
On the rotationally consistent triangle sector
\[
R_{xz}=R_{xy}R_{yz},
\]
the closed affine loop has
\[
\boxed{R_C=I},
\qquad
\boxed{
\mathbf t_C
=\operatorname{vec}(\mathcal T_{xyz})
=-\operatorname{vec}(\mathcal C_{xyz}).
}
\]
The scalar loop witness is
\[
\boxed{
\tau_C
=\|\mathbf t_C\|
=\sqrt{\frac12\operatorname{Tr}(\mathcal T_{xyz}^2)}.
}
\]
The source theorem is
\path{TIR/foundations/TIR_UNIVERSAL_LOOP_TORSION_SOURCE_BINDING_V0_1.md}.
\vTwelveStatus{B}{--}{PASS}

\section{Gate A2: Cartan refinement}

For an admitted shape-regular smooth family of shrinking relational triangles,
with edge size \(\ell\to0\), the covariant edge realization gives
\[
\mathcal E_{pq}
=e^a{}_{\mu}(p)\Delta x^{\mu}\sigma_a+O(\ell^2).
\]
With oriented area bivector \(\Sigma^{\mu\nu}_{\triangle}=O(\ell^2)\),
the discrete solder object satisfies
\[
\boxed{
\mathcal T_{\triangle}
=\frac12T^a{}_{\mu\nu}(x)
\Sigma^{\mu\nu}_{\triangle}\sigma_a
+O(\ell^3),
}
\]
where
\[
\boxed{T^a=de^a+\omega^a{}_b\wedge e^b}.
\]
The rotational loop satisfies
\[
\boxed{
R_{\triangle}
=I+\frac12\Omega_{\mu\nu}(x)
\Sigma^{\mu\nu}_{\triangle}+O(\ell^3),
}
\]
with
\[
\boxed{
\Omega^a{}_b
=d\omega^a{}_b+\omega^a{}_c\wedge\omega^c{}_b.
}
\]
Thus the corresponding area-normalized quantities converge to the Cartan
torsion and curvature components under the declared refinement assumptions.
The full affine-loop translation differs from the discrete solder torsion only
by \(O(\ell^3)\) on a smooth curved connection, so its area-normalized limit
carries the same torsion coefficient.

Canonical source:
\path{TIR/foundations/TIR_CARTAN_CONTINUUM_REFINEMENT_V0_1.md}.
\vTwelveStatus{B}{--}{PASS}

This is a conditional local/refining theorem.  It does not assert that the
actual global TIR relational complex already supplies the required production
refining family.

\section{Gate A3: zero torsion and Levi-Civita selection}

When the direct \(x\to z\) relation and the connection-composed
\(x\to y\to z\) description are admitted as two representations of the same
primitive ordered endpoint relation in the same comparison frame, intrinsic
affine-displacement uniqueness selects
\[
\boxed{\mathcal C_{xyz}=0}.
\]
Hence
\[
\mathcal T_{xyz}=0,
\]
and Gate A2 carries the endpoint-compatible regular refinement to
\[
\boxed{T^a=0}.
\]
Because the spatial transport is induced through
\[
SU(2)\xrightarrow{\operatorname{Ad}}SO(3),
\]
it preserves the spatial metric, so \(Dh=0\).  The standard uniqueness theorem
for a torsion-free metric-compatible connection then gives
\[
\boxed{D=D^{\rm LC}}.
\]
Zero torsion does not force zero curvature; the sector
\[
T^a=0,
\qquad
\Omega^a{}_b\ne0
\]
remains admitted.

Canonical source:
\path{TIR/foundations/TIR_ZERO_TORSION_LEVI_CIVITA_SELECTION_V0_1.md}.
\vTwelveStatus{B}{--}{PASS}

The broader context-lifted torsional Cartan branch remains typed separately.

\section{Gate A4: leading-loop locality and metric jet}

For a regular shrinking loop with \(A_C=\Theta(\epsilon^2)\), the leading
rotational coefficient is
\[
\boxed{
\frac{R_C-I}{A_C}\longrightarrow\Omega.
}
\]
Curvature-derivative terms enter at higher powers of the refinement scale.
The TIR Leading Refinement Rule selects the lowest finite nonzero covariant
normalized loop coefficient as the primitive continuum carrier.  On the A3
Levi-Civita sector, this gives the leading metric operator class
\[
\boxed{
\mathcal E_{\mu\nu}
=\mathcal E_{\mu\nu}(g,\partial g,\partial^2g),
}
\]
while \(\nabla\Omega,\nabla^2\Omega,\ldots\) remain typed as higher-jet
correction sectors.

Canonical source:
\path{TIR/foundations/TIR_LEADING_LOOP_LOCALITY_METRIC_JET_V0_1.md}.
\vTwelveStatus{B}{--}{PASS}

\section{Gate A5: global three-manifold certificate}

For a supplied finite pure tetrahedral complex, Gate A5 checks the closed
combinatorial three-manifold conditions through tetrahedral incidence and
spherical vertex links.  A passing certificate gives a combinatorial
three-manifold, and the standard Moise bridge supplies compatible smooth
realization in dimension three.  On such a carrier, compatible local
\(SO(3)\)-related coframes glue the metric globally, and the local A3
Levi-Civita restrictions glue by uniqueness.

Canonical source:
\path{TIR/foundations/TIR_GLOBAL_3MANIFOLD_SMOOTH_CERTIFICATE_V0_1.md}.

The certifier and theorem are implemented.  The actual source-owned TIR global
incidence complex has not yet been supplied as a frozen production input.
\vTwelveStatus{B}{--}{OPEN}

The production-input/freeze contract is
\path{TIR/foundations/TIR_GLOBAL_SPATIAL_COMPLEX_INPUT_CONTRACT_V0_1.md}.
The correct firewall is therefore
\[
\boxed{
\text{A5 certifier PASS}
\not\Rightarrow
\text{production global-complex PASS}.
}
\]

\section{Current spatial dependency line}

The synchronized spatial chain is
\[
\boxed{
\begin{aligned}
\mathbb C^2
&\to\operatorname{Herm}_0(2)\cong\mathbb R^3
\to\mathcal E_{xy}
\to W_{ij}^{X}/R_{ij}\\
&\to\mathcal T_{xyz}
\to\text{A2 Cartan refinement}
\to\text{A3 }T^a=0/D^{LC}\\
&\to\text{A4 leading metric jet}
\to\text{A5 global manifold certificate}.
\end{aligned}
}
\]
The remaining spatial/global gate is the frozen production incidence carrier,
followed by the inter-leaf temporal/spacetime join described in the completion
frontier.

\section{Source and validation receipt}

The canonical source chain now includes
\begin{itemize}
\item \path{TIR/foundations/TIR_WIJ_HOLONOMY_CROSSWALK_V0_1.md};
\item \path{TIR/foundations/TIR_SE3_ANCHOR_SOURCE_BINDING_V0_1.md};
\item \path{TIR/foundations/TIR_DISCRETE_SOLDER_FORM_V0_1.md};
\item \path{TIR/foundations/TIR_UNIVERSAL_LOOP_TORSION_SOURCE_BINDING_V0_1.md};
\item \path{TIR/foundations/TIR_CARTAN_CONTINUUM_REFINEMENT_V0_1.md};
\item \path{TIR/foundations/TIR_ZERO_TORSION_LEVI_CIVITA_SELECTION_V0_1.md};
\item \path{TIR/foundations/TIR_LEADING_LOOP_LOCALITY_METRIC_JET_V0_1.md};
\item \path{TIR/foundations/TIR_GLOBAL_3MANIFOLD_SMOOTH_CERTIFICATE_V0_1.md};
\item \path{TIR/foundations/TIR_GLOBAL_SPATIAL_COMPLEX_INPUT_CONTRACT_V0_1.md}.
\end{itemize}
The original deterministic v12 receipt remains generated by
\path{TIR/validation/tir_v12_ch07_holonomy_se3_torsion_v0_1.py}; the v12.1
repository-synchronization gate additionally checks the A2--A5 ownership
markers.

% ---- END TIR/monograph/v12/chapters/ch07_connections_holonomy_se3_solder_torsion.tex ----
\clearpage
